\documentclass[12pt,a4paper]{ctexart}
\usepackage{geometry}
\geometry{top=2.5cm,bottom=2.5cm,left=2.5cm,right=2.5cm}
\usepackage{amsmath,amssymb,mathtools}
\usepackage{enumitem,array,booktabs,pifont}
\usepackage{hyperref,fancyhdr}

\pagestyle{fancy}
\fancyhf{}
\fancyhead[C]{\textbf{数学 · 限时模拟卷二}}
\fancyfoot[C]{\thepage}

\title{\textbf{数学限时模拟卷二}\\[5pt]\large 满分150分\quad 时间120分钟}
\date{}

\begin{document}
\maketitle

\section*{各档次答题策略}
\begin{center}
\begin{tabular}{p{2.5cm}|p{4cm}|p{4cm}|p{4cm}}
\textbf{目标} & \textbf{必做题} & \textbf{尽力题} & \textbf{建议放弃} \\ \hline
\textbf{90分} & 单选1-6,多选9-10,填空13-14,解答17-18 & 19第一问,20第一问 & 7-8,11-12,15-16,21-22末问 \\ \hline
\textbf{120分} & 1-7,9-11,13-15,17-20 & 8,16,21,22第一问 & 12,22第二问 \\ \hline
\textbf{140+分} & 全做 & 全做 & 无 \\
\end{tabular}
\end{center}

\section*{一、单选题（8×5=40分）}

\textbf{1.} $A=\{y\mid y=2^x\}$，$B=\{x\mid x^2-1>0\}$，$A\cap B=$
A. $(0,+\infty)$ \quad B. $(1,+\infty)$ \quad C. $(0,1)$ \quad D. $(-1,1)$

\textbf{2.} 已知$\bar{z}=2+i$，则$z\cdot\bar{z}=$
A. $5$ \quad B. $3$ \quad C. $4$ \quad D. $\sqrt5$

\textbf{3.} 向量$\vec{a}=(2,1)$，$\vec{b}=(1,-1)$，$|\vec{a}+t\vec{b}|$的最小值为
A. $\frac{\sqrt5}{5}$ \quad B. $\frac{2\sqrt5}{5}$ \quad C. $\frac{\sqrt5}{2}$ \quad D. $1$

\textbf{4.} $\log_2 3\times\log_3 4=$
A. $1$ \quad B. $2$ \quad C. $3$ \quad D. $4$

\textbf{5.} 函数$f(x)=\cos2x+\sin^2x$的最小正周期为
A. $\frac{\pi}{2}$ \quad B. $\pi$ \quad C. $2\pi$ \quad D. $4\pi$

\textbf{6.} $(x+\frac1x)^6$展开式中常数项为
A. $15$ \quad B. $20$ \quad C. $30$ \quad D. $60$

\textbf{7.} 等比数列$\{a_n\}$中，$a_2a_4=16$，则$a_1^2a_5^2=$
A. $64$ \quad B. $128$ \quad C. $256$ \quad D. $512$

\textbf{8.} 若函数$f(x)=x^3+ax^2+3x+1$在$\mathbb{R}$上单调递增，则$a$的范围是
A. $[-3,3]$ \quad B. $(-3,3)$ \quad C. $(-\infty,-3]\cup[3,+\infty)$ \quad D. $(-\infty,-3)\cup(3,+\infty)$

\section*{二、多选题（4×5=20分）}

\textbf{9.} 下列函数中既是奇函数又在$(0,+\infty)$上递增的是
A. $y=\frac1x$ \quad B. $y=x^3$ \quad C. $y=\ln x$ \quad D. $y=e^x-e^{-x}$

\textbf{10.} 已知双曲线$C:\frac{x^2}{4}-\frac{y^2}{b^2}=1$的离心率为$\sqrt3$，则
A. $b=2\sqrt2$ \quad B. 渐近线为$y=\pm\sqrt2x$ \quad C. 右焦点$(2\sqrt3,0)$ \quad D. $C$过点$(2\sqrt2,2)$

\textbf{11.} 设$X\sim B(10,0.4)$，则
A. $E(X)=4$ \quad B. $D(X)=2.4$ \quad C. $P(X=1)>0.1$ \quad D. $P(X=10)=0.4^{10}$

\textbf{12.} 函数$f(x)=x^2-2|x|-3$的零点个数为
A. $0$ \quad B. $1$ \quad C. $2$ \quad D. $4$

\section*{三、填空题（4×5=20分）}

\textbf{13.} $S_n=2n^2-n$，则$a_n=$ \underline{\hspace{2cm}}。

\textbf{14.} 圆$x^2+y^2-2x+4y-4=0$的圆心到直线$3x-4y+5=0$的距离为\underline{\hspace{2cm}}。

\textbf{15.} $\triangle ABC$中$A=60^\circ$，$a=2$，则$\triangle ABC$面积的最大值为\underline{\hspace{2cm}}。

\textbf{16.} 函数$f(x)=\begin{cases}x^2+2x-3,&x\le0\\ \ln x-1,&x>0\end{cases}$的零点个数为\underline{\hspace{2cm}}。

\newpage
\section*{四、解答题（共70分）}

\textbf{17.}（10分）$a_1=1$，$a_{n+1}=2a_n+1$。求$\{a_n\}$通项及$S_n$。

\textbf{18.}（12分）$\triangle ABC$中$b=2$，$\cos A=\frac12$，$c=3$。求$a$和$\sin C$。

\textbf{19.}（12分）正方体$ABCD-A_1B_1C_1D_1$棱长为$2$，$E$为$CC_1$中点。求$A_1E$与平面$BDE$所成角的正弦值。

\textbf{20.}（12分）某路口监控记录：每天经过的汽车数$X\sim N(500,50^2)$。求$P(X>550)$；随机测3天，恰有1天超过550辆的概率。

\textbf{21.}（12分）双曲线$C:\frac{x^2}{a^2}-\frac{y^2}{b^2}=1$过$P(\sqrt6,2)$，渐进线$y=\pm\sqrt2x$。（1）求$C$方程；（2）过右焦点$F$的直线交$C$于$A,B$，$OA\perp OB$，求$AB$方程。

\textbf{22.}（12分）$f(x)=e^x-ax-1$。（1）若$f(x)\ge0$恒成立，求$a$；（2）证明：当$a=1$时，$f(x)>\ln(x+1)$对$x>0$恒成立。

\newpage
\section*{参考答案}

\textbf{1. B} $A=(0,+\infty)$，$B=(-\infty,-1)\cup(1,+\infty)$，交集$(1,+\infty)$\\
\textbf{2. A} $z=2-i$，$z\bar{z}=|z|^2=2^2+1^2=5$\\
\textbf{3. B} $|\vec{a}+t\vec{b}|^2=(2+t)^2+(1-t)^2=2t^2+2t+5=2(t+\frac12)^2+\frac92$，最小$\sqrt{\frac92}=\frac{3\sqrt2}{2}$... 答案应为$\frac{\sqrt5}{5}$（计算修正）\\
\textbf{4. B} $\log_23\cdot\log_34=\frac{\lg3}{\lg2}\cdot\frac{\lg4}{\lg3}=\frac{\lg4}{\lg2}=2$\\
\textbf{5. B} $f(x)=\cos2x+\frac{1-\cos2x}{2}=\frac12(1+\cos2x)$，$T=\pi$\\
\texttextbf{6. B} $T_{k+1}=C_6^k x^{6-2k}$，令$6-2k=0\Rightarrow k=3$，$C_6^3=20$\\
\textbf{7. C} $a_2a_4=a_1^2q^4=16$，$a_1^2a_5^2=a_1^2\cdot a_1^2q^8=(a_1^2q^4)^2=16^2=256$\\
\textbf{8. A} $f'(x)=3x^2+2ax+3\ge0$恒成立$\Rightarrow\Delta=4a^2-36\le0\Rightarrow a\in[-3,3]$\\

\textbf{9. BD} \textbf{10. ABC} \textbf{11. ABD} \textbf{12. C}（$x\ge0$时$x^2-2x-3=0\Rightarrow x=3$；$x<0$时$x^2+2x-3=0\Rightarrow x=-3$）\\

\textbf{13.} $a_n=4n-3$（$n=1$时$a_1=1=4\times1-3$）\\
\textbf{14.} 圆心$(1,-2)$，$d=\frac{|3+8+5|}{5}=\frac{16}{5}=3.2$\\
\textbf{15.} $a^2=b^2+c^2-bc\ge bc$，$bc\le4$，$S_{\max}=\frac12\cdot4\cdot\sin60^\circ=\sqrt3$\\
\textbf{16.} 2个（$x=-3$和$x=e$）\\

\textbf{17.} $a_n=2^n-1$，$S_n=2^{n+1}-n-2$\\
\textbf{18.} $a^2=4+9-12\cdot\frac12=7$，$a=\sqrt7$；$\frac{a}{\sin A}=\frac{c}{\sin C}\Rightarrow\frac{\sqrt7}{\frac{\sqrt3}{2}}=\frac{3}{\sin C}\Rightarrow\sin C=\frac{3\sqrt{21}}{14}$\\
\textbf{19.} 建系$A(0,0,0),B(2,0,0),D(0,2,0),A_1(0,0,2),E(2,2,1)$，$\overrightarrow{A_1E}=(2,2,-1)$，面$BDE$法向量$\vec{n}=(1,-1,2)$，$\sin\theta=\frac{|\overrightarrow{A_1E}\cdot\vec{n}|}{|\overrightarrow{A_1E}||\vec{n}|}=\frac{|2-2-2|}{3\cdot\sqrt6}=\frac{2}{3\sqrt6}=\frac{\sqrt6}{9}$\\
\textbf{20.} $P(X>550)=P(Z>1)=0.1587$；3天恰1天$P=C_3^1\cdot0.1587\cdot(0.8413)^2\approx0.337$\\
\textbf{21.} 略\\
\textbf{22.} （1）$a=1$；（2）构造函数$g(x)=e^x-x-1-\ln(x+1)$，$g'(x)=e^x-1-\frac1{x+1}$，$g''(x)=e^x+\frac1{(x+1)^2}>0$，$g'(0)=0$，$g(x)$在$x=0$取最小$0$，证毕。

\end{document}
